\IfFileExists{stacks-project.cls}{%
\documentclass{stacks-project}
}{%
\documentclass{amsart}
}

% For dealing with references we use the comment environment
\usepackage{verbatim}
\newenvironment{reference}{\comment}{\endcomment}
%\newenvironment{reference}{}{}
\newenvironment{slogan}{\comment}{\endcomment}
\newenvironment{history}{\comment}{\endcomment}

% For commutative diagrams we use Xy-pic
\usepackage[all]{xy}

% We use 2cell for 2-commutative diagrams.
\xyoption{2cell}
\UseAllTwocells

% We use multicol for the list of chapters between chapters
\usepackage{multicol}

% This is generally recommended for better output
\usepackage{lmodern}
\usepackage[T1]{fontenc}

% For cross-file-references
\usepackage{xr-hyper}

% Package for hypertext links:
\usepackage{hyperref}

% For any local file, say "hello.tex" you want to link to please
% use \externaldocument[hello-]{hello}
\externaldocument[introduction-]{introduction}
\externaldocument[conventions-]{conventions}
\externaldocument[sets-]{sets}
\externaldocument[categories-]{categories}
\externaldocument[topology-]{topology}
\externaldocument[sheaves-]{sheaves}
\externaldocument[sites-]{sites}
\externaldocument[stacks-]{stacks}
\externaldocument[fields-]{fields}
\externaldocument[algebra-]{algebra}
\externaldocument[brauer-]{brauer}
\externaldocument[homology-]{homology}
\externaldocument[derived-]{derived}
\externaldocument[simplicial-]{simplicial}
\externaldocument[more-algebra-]{more-algebra}
\externaldocument[smoothing-]{smoothing}
\externaldocument[modules-]{modules}
\externaldocument[sites-modules-]{sites-modules}
\externaldocument[injectives-]{injectives}
\externaldocument[cohomology-]{cohomology}
\externaldocument[sites-cohomology-]{sites-cohomology}
\externaldocument[dga-]{dga}
\externaldocument[dpa-]{dpa}
\externaldocument[sdga-]{sdga}
\externaldocument[hypercovering-]{hypercovering}
\externaldocument[schemes-]{schemes}
\externaldocument[constructions-]{constructions}
\externaldocument[properties-]{properties}
\externaldocument[morphisms-]{morphisms}
\externaldocument[coherent-]{coherent}
\externaldocument[divisors-]{divisors}
\externaldocument[limits-]{limits}
\externaldocument[varieties-]{varieties}
\externaldocument[topologies-]{topologies}
\externaldocument[descent-]{descent}
\externaldocument[perfect-]{perfect}
\externaldocument[more-morphisms-]{more-morphisms}
\externaldocument[flat-]{flat}
\externaldocument[groupoids-]{groupoids}
\externaldocument[more-groupoids-]{more-groupoids}
\externaldocument[etale-]{etale}
\externaldocument[chow-]{chow}
\externaldocument[intersection-]{intersection}
\externaldocument[pic-]{pic}
\externaldocument[weil-]{weil}
\externaldocument[adequate-]{adequate}
\externaldocument[dualizing-]{dualizing}
\externaldocument[duality-]{duality}
\externaldocument[discriminant-]{discriminant}
\externaldocument[derham-]{derham}
\externaldocument[local-cohomology-]{local-cohomology}
\externaldocument[algebraization-]{algebraization}
\externaldocument[curves-]{curves}
\externaldocument[resolve-]{resolve}
\externaldocument[models-]{models}
\externaldocument[functors-]{functors}
\externaldocument[equiv-]{equiv}
\externaldocument[pione-]{pione}
\externaldocument[etale-cohomology-]{etale-cohomology}
\externaldocument[proetale-]{proetale}
\externaldocument[relative-cycles-]{relative-cycles}
\externaldocument[more-etale-]{more-etale}
\externaldocument[trace-]{trace}
\externaldocument[crystalline-]{crystalline}
\externaldocument[spaces-]{spaces}
\externaldocument[spaces-properties-]{spaces-properties}
\externaldocument[spaces-morphisms-]{spaces-morphisms}
\externaldocument[decent-spaces-]{decent-spaces}
\externaldocument[spaces-cohomology-]{spaces-cohomology}
\externaldocument[spaces-limits-]{spaces-limits}
\externaldocument[spaces-divisors-]{spaces-divisors}
\externaldocument[spaces-over-fields-]{spaces-over-fields}
\externaldocument[spaces-topologies-]{spaces-topologies}
\externaldocument[spaces-descent-]{spaces-descent}
\externaldocument[spaces-perfect-]{spaces-perfect}
\externaldocument[spaces-more-morphisms-]{spaces-more-morphisms}
\externaldocument[spaces-flat-]{spaces-flat}
\externaldocument[spaces-groupoids-]{spaces-groupoids}
\externaldocument[spaces-more-groupoids-]{spaces-more-groupoids}
\externaldocument[bootstrap-]{bootstrap}
\externaldocument[spaces-pushouts-]{spaces-pushouts}
\externaldocument[spaces-chow-]{spaces-chow}
\externaldocument[groupoids-quotients-]{groupoids-quotients}
\externaldocument[spaces-more-cohomology-]{spaces-more-cohomology}
\externaldocument[spaces-simplicial-]{spaces-simplicial}
\externaldocument[spaces-duality-]{spaces-duality}
\externaldocument[formal-spaces-]{formal-spaces}
\externaldocument[restricted-]{restricted}
\externaldocument[spaces-resolve-]{spaces-resolve}
\externaldocument[formal-defos-]{formal-defos}
\externaldocument[defos-]{defos}
\externaldocument[cotangent-]{cotangent}
\externaldocument[examples-defos-]{examples-defos}
\externaldocument[algebraic-]{algebraic}
\externaldocument[examples-stacks-]{examples-stacks}
\externaldocument[stacks-sheaves-]{stacks-sheaves}
\externaldocument[criteria-]{criteria}
\externaldocument[artin-]{artin}
\externaldocument[quot-]{quot}
\externaldocument[stacks-properties-]{stacks-properties}
\externaldocument[stacks-morphisms-]{stacks-morphisms}
\externaldocument[stacks-limits-]{stacks-limits}
\externaldocument[stacks-cohomology-]{stacks-cohomology}
\externaldocument[stacks-perfect-]{stacks-perfect}
\externaldocument[stacks-introduction-]{stacks-introduction}
\externaldocument[stacks-more-morphisms-]{stacks-more-morphisms}
\externaldocument[stacks-geometry-]{stacks-geometry}
\externaldocument[moduli-]{moduli}
\externaldocument[moduli-curves-]{moduli-curves}
\externaldocument[examples-]{examples}
\externaldocument[exercises-]{exercises}
\externaldocument[guide-]{guide}
\externaldocument[desirables-]{desirables}
\externaldocument[coding-]{coding}
\externaldocument[obsolete-]{obsolete}
\externaldocument[fdl-]{fdl}
\externaldocument[index-]{index}

% Theorem environments.
%
\theoremstyle{plain}
\newtheorem{theorem}[subsection]{Theorem}
\newtheorem{proposition}[subsection]{Proposition}
\newtheorem{lemma}[subsection]{Lemma}

\theoremstyle{definition}
\newtheorem{definition}[subsection]{Definition}
\newtheorem{example}[subsection]{Example}
\newtheorem{exercise}[subsection]{Exercise}
\newtheorem{situation}[subsection]{Situation}

\theoremstyle{remark}
\newtheorem{remark}[subsection]{Remark}
\newtheorem{remarks}[subsection]{Remarks}

\numberwithin{equation}{subsection}

% Macros
%
\def\lim{\mathop{\mathrm{lim}}\nolimits}
\def\colim{\mathop{\mathrm{colim}}\nolimits}
\def\Spec{\mathop{\mathrm{Spec}}}
\def\Hom{\mathop{\mathrm{Hom}}\nolimits}
\def\Ext{\mathop{\mathrm{Ext}}\nolimits}
\def\SheafHom{\mathop{\mathcal{H}\!\mathit{om}}\nolimits}
\def\SheafExt{\mathop{\mathcal{E}\!\mathit{xt}}\nolimits}
\def\Sch{\mathit{Sch}}
\def\Mor{\mathop{\mathrm{Mor}}\nolimits}
\def\Ob{\mathop{\mathrm{Ob}}\nolimits}
\def\Sh{\mathop{\mathit{Sh}}\nolimits}
\def\NL{\mathop{N\!L}\nolimits}
\def\CH{\mathop{\mathrm{CH}}\nolimits}
\def\proetale{{pro\text{-}\acute{e}tale}}
\def\etale{{\acute{e}tale}}
\def\QCoh{\mathit{QCoh}}
\def\Ker{\mathop{\mathrm{Ker}}}
\def\Im{\mathop{\mathrm{Im}}}
\def\Coker{\mathop{\mathrm{Coker}}}
\def\Coim{\mathop{\mathrm{Coim}}}

% Boxtimes
%
\DeclareMathSymbol{\boxtimes}{\mathbin}{AMSa}{"02}

%
% Macros for moduli stacks/spaces
%
\def\QCohstack{\mathcal{QC}\!\mathit{oh}}
\def\Cohstack{\mathcal{C}\!\mathit{oh}}
\def\Spacesstack{\mathcal{S}\!\mathit{paces}}
\def\Quotfunctor{\mathrm{Quot}}
\def\Hilbfunctor{\mathrm{Hilb}}
\def\Curvesstack{\mathcal{C}\!\mathit{urves}}
\def\Polarizedstack{\mathcal{P}\!\mathit{olarized}}
\def\Complexesstack{\mathcal{C}\!\mathit{omplexes}}
% \Pic is the operator that assigns to X its picard group, usage \Pic(X)
% \Picardstack_{X/B} denotes the Picard stack of X over B
% \Picardfunctor_{X/B} denotes the Picard functor of X over B
\def\Pic{\mathop{\mathrm{Pic}}\nolimits}
\def\Picardstack{\mathcal{P}\!\mathit{ic}}
\def\Picardfunctor{\mathrm{Pic}}
\def\Deformationcategory{\mathcal{D}\!\mathit{ef}}
\usepackage{cleveref}
\crefname{originallemma}{original lemma}{original lemmas}
\Crefname{originallemma}{Original lemma}{Original lemmas}
\makeatletter
\newlabel{varieties-lemma-geometrically-connected-if-connected-and-point}{{04KV}{source}{}{sourcepointer-1}{}}
\newlabel{varieties-lemma-geometrically-connected-if-connected-and-point@cref}{{[originallemma][0][]04KV}{[1][0][]source}{}{}{}}
\makeatother
\crefname{originallemma}{original lemma}{original lemmas}
\Crefname{originallemma}{Original lemma}{Original lemmas}
\makeatletter
\newlabel{varieties-lemma-make-Jacobson}{{0479}{source}{}{sourcepointer-2}{}}
\newlabel{varieties-lemma-make-Jacobson@cref}{{[originallemma][0][]0479}{[1][0][]source}{}{}{}}
\makeatother
\crefname{originallemma}{original lemma}{original lemmas}
\Crefname{originallemma}{Original lemma}{Original lemmas}
\makeatletter
\newlabel{lemma-group-scheme-field-geometrically-irreducible}{{047M}{source}{}{sourcepointer-3}{}}
\newlabel{lemma-group-scheme-field-geometrically-irreducible@cref}{{[originallemma][0][]047M}{[1][0][]source}{}{}{}}
\makeatother
\crefname{originallemma}{original lemma}{original lemmas}
\Crefname{originallemma}{Original lemma}{Original lemmas}
\makeatletter
\newlabel{varieties-lemma-geometrically-reduced-any-base-change}{{035Z}{source}{}{sourcepointer-4}{}}
\newlabel{varieties-lemma-geometrically-reduced-any-base-change@cref}{{[originallemma][0][]035Z}{[1][0][]source}{}{}{}}
\makeatother
\crefname{originallemma}{original lemma}{original lemmas}
\Crefname{originallemma}{Original lemma}{Original lemmas}
\makeatletter
\newlabel{varieties-lemma-bijection-irreducible-components}{{038F}{source}{}{sourcepointer-5}{}}
\newlabel{varieties-lemma-bijection-irreducible-components@cref}{{[originallemma][0][]038F}{[1][0][]source}{}{}{}}
\makeatother
\crefname{originallemma}{original lemma}{original lemmas}
\Crefname{originallemma}{Original lemma}{Original lemmas}
\makeatletter
\newlabel{lemma-irreducible-group-scheme-over-field-qc}{{0B7P}{source}{}{sourcepointer-6}{}}
\newlabel{lemma-irreducible-group-scheme-over-field-qc@cref}{{[originallemma][0][]0B7P}{[1][0][]source}{}{}{}}
\makeatother
\crefname{originallemma}{original lemma}{original lemmas}
\Crefname{originallemma}{Original lemma}{Original lemmas}
\makeatletter
\newlabel{lemma-group-scheme-over-field-translate-open}{{0B7N}{source}{}{sourcepointer-7}{}}
\newlabel{lemma-group-scheme-over-field-translate-open@cref}{{[originallemma][0][]0B7N}{[1][0][]source}{}{}{}}
\makeatother
\crefname{originallemma}{original lemma}{original lemmas}
\Crefname{originallemma}{Original lemma}{Original lemmas}
\makeatletter
\newlabel{lemma-group-scheme-over-field-separated}{{047L}{source}{}{sourcepointer-8}{}}
\newlabel{lemma-group-scheme-over-field-separated@cref}{{[originallemma][0][]047L}{[1][0][]source}{}{}{}}
\makeatother
\crefname{originallemma}{original lemma}{original lemmas}
\Crefname{originallemma}{Original lemma}{Original lemmas}
\makeatletter
\newlabel{schemes-lemma-quasi-compact-permanence}{{03GI}{source}{}{sourcepointer-9}{}}
\newlabel{schemes-lemma-quasi-compact-permanence@cref}{{[originallemma][0][]03GI}{[1][0][]source}{}{}{}}
\makeatother
\crefname{originallemma}{original lemma}{original lemmas}
\Crefname{originallemma}{Original lemma}{Original lemmas}
\makeatletter
\newlabel{morphisms-lemma-reach-points-scheme-theoretic-image}{{02JQ}{source}{}{sourcepointer-10}{}}
\newlabel{morphisms-lemma-reach-points-scheme-theoretic-image@cref}{{[originallemma][0][]02JQ}{[1][0][]source}{}{}{}}
\makeatother
\title{A connected group scheme over a field is irreducible}
\author{The Stacks Project Authors}
\date{stacks-project:groupoids.tex:local-snapshot-20261012; GNU FDL 1.2 or later}
\begin{document}
\maketitle
\setcounter{subsection}{1}
\noindent\textit{Editorial note.} All delivered mathematical text is the author's own, copied from the pinned Stacks Project chapter file groupoids.tex; no statement or proof step was rewritten, reordered or removed. The counted span is the author environment and its own proof environment. References inside the span are handed to the assembler as navigation stubs that carry the canonical tag of the cited tag only; whether each of them is load-bearing is recorded on the stub itself.\par\medskip
\noindent\textit{Allowed standard prerequisites (editorial).}\begin{itemize}
\item Lemma varieties-lemma-geometrically-connected-if-connected-and-point (canonical tag 04KV). cited at proof lines 1172-1174. G is geometrically connected, so irreducibility can be checked after a field extension.
\item Lemma varieties-lemma-make-Jacobson (canonical tag 0479). cited at proof lines 1176-1179. Reduces to k algebraically closed, G Jacobson, closed points k-rational.
\item Lemma lemma-group-scheme-field-geometrically-irreducible (canonical tag 047M). cited at proof lines 1182-1185 and 1196-1198. The component Z through e exists and distinct components are disjoint.
\item Varieties lemmas 035Z and 038F. cited at proof lines 1186-1189. Z x\_k Z is reduced and irreducible, so multiplication factors through Z.
\item Lemma lemma-irreducible-group-scheme-over-field-qc (0B7P) and lemma-group-scheme-over-field-translate-open (0B7N). cited at proof lines 1197-1203. Z quasi-compact and W open.
\item Lemma lemma-group-scheme-over-field-separated (047L), schemes-lemma-quasi-compact-permanence (03GI), morphisms-lemma-reach-points-scheme-theoretic-image (02JQ). cited at proof lines 1209-1216. The specialization argument proving W is closed.
\end{itemize}\medskip
\ifdefined\hypertarget\hypertarget{sourcepointer-1}{}\fi
\noindent\textit{Source reference.} Original lemma 04KV: sources/stacks-chow/tags.txt, source lines 5986--5986. Cited at proof lines 1172-1174 to see G is geometrically connected; load-bearing for the reduction.\par\medskip
\ifdefined\hypertarget\hypertarget{sourcepointer-2}{}\fi
\noindent\textit{Source reference.} Original lemma 0479: sources/stacks-chow/tags.txt, source lines 5500--5500. Cited at proof lines 1176-1179 to reduce to k algebraically closed, G Jacobson with rational closed points; load-bearing.\par\medskip
\ifdefined\hypertarget\hypertarget{sourcepointer-3}{}\fi
\noindent\textit{Source reference.} Original lemma 047M: sources/stacks-chow/tags.txt, source lines 5513--5513. Cited at proof lines 1182-1185 (existence of the component Z through e, disjointness of distinct components); load-bearing.\par\medskip
\ifdefined\hypertarget\hypertarget{sourcepointer-4}{}\fi
\noindent\textit{Source reference.} Original lemma 035Z: sources/stacks-chow/tags.txt, source lines 4146--4146. Cited at proof lines 1186-1189 to see Z x\_k Z is reduced and irreducible; load-bearing for Z being a subgroup scheme.\par\medskip
\ifdefined\hypertarget\hypertarget{sourcepointer-5}{}\fi
\noindent\textit{Source reference.} Original lemma 038F: sources/stacks-chow/tags.txt, source lines 4234--4234. Cited at proof lines 1186-1189 (with 035Z) to see Z x\_k Z is irreducible; load-bearing.\par\medskip
\ifdefined\hypertarget\hypertarget{sourcepointer-6}{}\fi
\noindent\textit{Source reference.} Original lemma 0B7P: sources/stacks-chow/tags.txt, source lines 14112--14112. Cited at proof lines 1197-1199 to know Z is quasi-compact; load-bearing for choosing U.\par\medskip
\ifdefined\hypertarget\hypertarget{sourcepointer-7}{}\fi
\noindent\textit{Source reference.} Original lemma 0B7N: sources/stacks-chow/tags.txt, source lines 14111--14111. Cited at proof lines 1201-1203 to see W = image of Z x\_k U -> G is open; load-bearing.\par\medskip
\ifdefined\hypertarget\hypertarget{sourcepointer-8}{}\fi
\noindent\textit{Source reference.} Original lemma 047L: sources/stacks-chow/tags.txt, source lines 983--983. Cited at proof lines 1209-1211 to run the quasi-compact permanence argument; load-bearing.\par\medskip
\ifdefined\hypertarget\hypertarget{sourcepointer-9}{}\fi
\noindent\textit{Source reference.} Original lemma 03GI: sources/stacks-chow/tags.txt, source lines 4529--4529. Cited at proof lines 1209-1212 to see W -> G is quasi-compact; load-bearing.\par\medskip
\ifdefined\hypertarget\hypertarget{sourcepointer-10}{}\fi
\noindent\textit{Source reference.} Original lemma 02JQ: sources/stacks-chow/tags.txt, source lines 3345--3345. Cited at proof lines 1213-1216 to reduce closedness of W to containing the irreducible components meeting it; load-bearing.\par\medskip
\begin{lemma}
\label{lemma-connected-group-scheme-over-field-irreducible}
Let $G$ be a group scheme over a field $k$. If $G$ is connected,
then $G$ is irreducible.
\end{lemma}

\begin{proof}
By Varieties, Lemma
\ref{varieties-lemma-geometrically-connected-if-connected-and-point}
we see that $G$ is geometrically connected. If we show that $G_K$
is irreducible for some field extension $K/k$, then
the lemma follows. Hence we may apply
Varieties, Lemma \ref{varieties-lemma-make-Jacobson}
to reduce to the case where $k$ is algebraically closed,
$G$ is a Jacobson scheme, and all the closed points are $k$-rational.

\medskip\noindent
Let $Z \subset G$ be the unique irreducible component of $G$ passing
through the neutral element, see
Lemma \ref{lemma-group-scheme-field-geometrically-irreducible}.
Endowing $Z$ with the reduced induced closed subscheme structure,
we see that $Z \times_k Z$ is reduced and irreducible
(Varieties, Lemmas
\ref{varieties-lemma-geometrically-reduced-any-base-change} and
\ref{varieties-lemma-bijection-irreducible-components}).
We conclude that $m|_{Z \times_k Z} : Z \times_k Z \to G$ factors
through $Z$. Hence $Z$ becomes a closed subgroup scheme of $G$.

\medskip\noindent
To get a contradiction, assume there exists another irreducible
component $Z' \subset G$. Then $Z \cap Z' = \emptyset$ by
Lemma \ref{lemma-group-scheme-field-geometrically-irreducible}.
By Lemma \ref{lemma-irreducible-group-scheme-over-field-qc}
we see that $Z$ is quasi-compact. Thus we may choose a quasi-compact open
$U \subset G$ with $Z \subset U$ and $U \cap Z' = \emptyset$.
The image $W$ of $Z \times_k U \to G$ is open in $G$ by
Lemma \ref{lemma-group-scheme-over-field-translate-open}.
On the other hand, $W$ is quasi-compact as the image of a
quasi-compact space. We claim that $W$ is closed.

\medskip\noindent
Proof of the claim. Since $W$ is quasi-compact and $G$ is separated
(Lemma \ref{lemma-group-scheme-over-field-separated}), the inclusion
$W \to G$ is quasi-compect by Schemes, Lemma
\ref{schemes-lemma-quasi-compact-permanence}. We conclude that
points in the closure of $W$ are specializations of points of $W$
(Morphisms, Lemma \ref{morphisms-lemma-reach-points-scheme-theoretic-image}).
Thus we have to show that any irreducible
component $Z'' \subset G$ of $G$ which meets $W$ is contained in $W$.
As $G$ is Jacobson and closed points are rational,
$Z'' \cap W$ has a rational point
$g \in Z''(k) \cap W(k)$ and hence $Z'' = Zg$. But $W = m(Z \times_k W)$
by construction, so $Z'' \cap W \not = \emptyset$ implies
$Z'' \subset W$.

\medskip\noindent
By the claim $W \subset G$ is an open and closed subset of $G$.
Now $W \cap Z' = \emptyset$ since otherwise by the argument given in
the preceding paragraph we would get $Z' = Zg$ for some $g \in W(k)$.
Then as $Z$ is a subgroup we could even pick $g \in U(k)$ which would
contradict $Z' \cap U = \emptyset$. Hence $W \subset G$ is a proper open
and closed subset which contradicts the assumption that $G$ is connected.
\end{proof}

\end{document}
